\begin{abstract}
Este artículo presenta un análisis de arquitecturas de antenas operando en bandas L y S para la gestión eficiente de redes de sensores ambientales mediante constelaciones de satélites en órbita baja (LEO). Se examinan requisitos de enlaces directos a sensores IoT de bajo consumo, considerando desafíos como Doppler, desvanecimiento, ganancia y polarización circular. Se comparan diseños conformes, phased-array y helicoidales desplegables, evaluando métricas de rendimiento como ganancia, eficiencia y cobertura global. Los resultados de simulaciones y análisis de estado del arte demuestran la viabilidad de soluciones híbridas para monitoreo en tiempo casi real de parámetros ambientales. Se discuten implicaciones para implementaciones CubeSat y estaciones terrestres, junto con recomendaciones para optimizaciones futuras en entornos dinámicos.
\end{abstract}